%\newpage
\section{Cauchy's 1815 Theorem}

Cauchy's 1815 theorem says that if a function of $n$ variables takes more than two distinct values when you permute the variables, then it must take at least $p$ distinct values, where $p$ is the largest prime less than or equal to $n$. 
The number of distinct values always divides $n$ factorial. Symmetric functions take one value, and the only way to get exactly two are by use of the alternating functions, like the Vandermonde product.

\begin{theorem}[Cauchy's theorem:] For a function of n variables, $f=f(x_1,\hdots,x_n)$, the number of distinct values under all permutations of the variables is either 1, 2, or at least $p$, where $p$ is the largest prime less than or equal to $n$. 
\end{theorem}
\begin{proof}
Let $f$ be a function of $n$ (distinct) variables $f=f(x_1,\hdots,x_n)$. 
Permuting the variables produces up to $n!$ possible inputs,  but we assume $f$ to take only $R$ distinct values under these permutations. 
By Lagrange's theorem  $n!$ is divisible $R$. Assume for contradiction purposes that $2 < R < p$. 
Consider a $p$-cycle $\sigma$ (a cyclic permutation of order $p$ on $p$ of the variables, fixing the rest).  
The sequence of function values obtained by applying successive powers of $\sigma$ to a fixed initial ntuple of variables:
\[ f(x), ~f(\sigma(x)), ~f(\sigma^2(x)), \hdots, ~f(\sigma^{p-1}(x)).\]
There are at most $p$ distinct values in the above list (since $\sigma^p$ is the identity).
Let \[ v_k = f(\sigma^k(x))  \mbox{ for } k = 0, 1, ..., p-1,\] stand for these $p$ values  and lest us apply the {\bf Pigeonhole Principle}
(i.e., if you have more pigeons than pigeonholes, at least one hole has more than one pigeon). 
These $p$ values are drawn from the set of $R$ assumed distinct possible values that $f$ can ever take, and we assume $R < p$. 
These values satisfy the cyclic relation: $v_{k+1} = f(\sigma^{k+1}(x))$ is obtained from $v_k$ by applying $\sigma \pmod{p}$.
By the pigeonhole principle at least two of these $p$ values must be equal.
\[ v_r = v_s \mbox{ for some } 0 \le s < r \le p-1. \]
Let $d = r - s$ (so $1 \le d \le p-1$).\\

Now apply powers of $\sigma$ repeatedly. 
Because $\sigma$ shifts the indices cyclically (i.e., applying $\sigma$ moves $v_k$ to $v_{(k+1)\pmod{p}}$),
the assumption $v_r = v_s$ means the sequence is periodic with a smaller step. 
The equality propagates under shifts:
Apply $\sigma$ any number of times to both sides of $v_r = v_s$.
This gives:
$v_{r + m} = v_{s + m \pmod{p}}$ for any integer $m$.
In other words, shifting the indices by the same amount preserves the equality.
The index differences are multiples of $d$:
The indices where we know equalities hold are exactly the index positions that differ by multiples of $d \pmod{p}$.
That is, for any integer $m$:
$v_k = v_{(k + m - d)  \pmod{p}}$, whenever the equality chain reaches there.\\

Primality of $p$ makes $d$ a generator:
Crucially because $p$ is prime and $d$ is between 1 and $p -1$, we have $\gcd(d, p) = 1$.
Therefore, $d$ has a multiplicative inverse modulo $p$, and the multiples
\[ \{ 0\cdot d, 1\cdot d, 2\cdot d,  \hdots,  (p-1)\cdot d \} \pmod{p} \]
hit every residue class modulo $p$ exactly once and we cycle through all residues modulo $p$. 
Starting from any fixed index $k$, we can reach any other index $l$ by adding a suitable multiple of $d$ modulo $p$.
Therefore $v_k = v_l \mbox{ for all } k, l$.
In group theory terms: the subgroup generated by $d$ modulo $p$ is the entire cyclic group $\Z/p\Z$. Therefore entire orbit of $f(\sigma^k x)$ for $k = 0, 1, 2, \hdots, p-1$  collapses to a single value. \\

{\bf Conclusion}:\\

The shifts by multiples of $d$ generate the entire cyclic group of order $p$.
It follows that all the $v_k$ are equal (the orbit collapses to a single value).
In particular, $v_0 = v_1$, i.e., $f(x) = f(\sigma(x))$.
Since the initial  ntuple $x$ was arbitrary, this holds for every starting assignment: $\sigma$ leaves $f$ invariant (does not change its value).
By choosing different $p$-cycles (or by conjugacy/relabeling of variables), this applies to all $p$-cycles in $S_n$.\\

This implies that all $p$-cycles leave $f$ invariant (fix every value).
One can then construct a 3-cycle as a product (or commutator) of two such $p$-cycles. Thus, 3-cycles also fix values.
A 3-cycle is a product of two transpositions. 
Analysis of transpositions then shows that either $f$ is fully symmetric $(R = 1)$ or it takes exactly two values (like the sign/alternating case, e.g., Vandermonde product).
But we assumed $R > 2$, a contradiction. You can't get three or more without hitting at least $p$ values.  \qed

\end{proof}


Thus, no such $R$ with $2 < R < p$ is possible. The pigeonhole step forces $p$-cycles (and then smaller cycles) to preserve values, collapsing the possibilities to $R = 1$ or $2$.
This is the heart of Cauchy's 1815 argument (sometimes phrased in terms of``forms'' or ``arrangement'' of the variables). 
Modern presentations often frame it in terms of the action of the symmetric group $S_n$ on the values, with stabilizers and orbits (see the {\bf Orbit-Stabilizer Theorem}) 
but the core pigeonhole idea on powers of a cycle remains the same.\\

{\bf Why Primality of p Is Essential}\\
If $p$ were composite, say $d$ shares a common factor with $p$, the multiples of $d \pmod{p}$ would only generate a proper subgroup, 
and you might only force equality within cosets rather than full invariance under $\sigma$. 
The primality ensures the generated subgroup is the whole group, forcing full invariance.
This is the precise bridge from the pigeonhole collision to invariance under the full $p$-cycle. 
Once you have invariance under all $p$-cycles, the rest of the proof proceeds by constructing smaller cycles (e.g., 3-cycles) as products/commutators and analyzing transpositions.\\
%\newpage
